\subsubsection{The improved critical exponent survives residual mixing}
\label{sec:residualfused}

The fixed potential of Appendix~\ref{app:fused} also controls
residual networks. The skip branch costs one child when it is
examined, but its effect on the potential has a useful concavity
term. That term prevents an additional factor from appearing in
the quadratic growth coefficient.

\begin{lemma}[Mixing a potential estimate with an identity step]
\label{lem:respotentialmix}
Let $0<\eta\le1/4$, $\phi(v)=\sqrt{1+\eta-v^2}$,
$u,w\in[-1,1]$, $0\le\omega\le1$, and $0\le\xi\le1$.
Then
\begin{equation}\label{eq:respotentialmix}
\log\frac{(1-\omega)\phi(u)+\omega e^\xi\phi(w)}
              {\phi((1-\omega)u+\omega w)}
\le\omega\xi+4\eta^{-3/2}\omega\xi^2.
\end{equation}
\end{lemma}
\paragraph{Proof idea.}
The concavity gap of the potential absorbs the mismatch between skip and
update means. A quadratic completion bounds the remaining product term
uniformly over the mixing probability.

\begin{proof}
Put $A=1+\eta$, $v=(1-\omega)u+\omega w$, and
$\Delta=|u-w|$. Direct differentiation gives
$|\phi'|\le\eta^{-1/2}$ and $-\phi''\ge A^{-1/2}$.
Thus the Jensen gap is at least
\[
\phi(v)-(1-\omega)\phi(u)-\omega\phi(w)
\ge\frac{\omega(1-\omega)}{2\sqrt A}\Delta^2.
\]
For example, this follows by applying concavity to
$\phi(x)+x^2/(2\sqrt A)$.
Also $|\phi(w)-\phi(v)|\le(1-\omega)\Delta/\sqrt\eta$.
Writing $t=e^\xi-1$, subtraction of
$(1+\omega t)\phi(v)$ from the numerator in
\eqref{eq:respotentialmix} therefore gives at most
\[
\omega(1-\omega)
 \left(\frac{t\Delta}{\sqrt\eta}
               -\frac{\Delta^2}{2\sqrt A}\right)
\le\frac{\omega t^2\sqrt A}{2\eta}.
\]
Divide by $\phi(v)\ge\sqrt\eta$.
For $0\le\xi\le1$, positivity of the exponential series gives
$e^\xi-1\le\xi+\xi^2\le2\xi$.
The resulting ratio is at most
$1+\omega\xi+4\eta^{-3/2}\omega\xi^2$.
Taking logarithms and using $\log(1+x)\le x$ proves the claim.
\end{proof}

For this variant, retain the radii from \eqref{eq:resrounded},
but use centers within $\alpha\varepsilon$ of the residual
midpoint $M$. Project these centers only onto $[-1,1]$.
The normalized interval invariant still holds, and the center
shift is now small enough for the potential estimate.
The root is implemented separately: draw the last active
nonlinear layer by a reverse geometric waiting time with parameter
$\alpha$. If there is no such layer, return $x_i$.
Otherwise use the raw tanh-row sampler at that layer.
This implements the exact identity
\begin{equation}\label{eq:resrawrootidentity}
h_{D,i}=(1-\alpha)^D x_i+
\alpha\sum_{j=1}^D(1-\alpha)^{D-j}
 \tanh\!\left(b_{j,i}+\sum_kw_{j,ik}h_{j-1,k}\right).
\end{equation}
All internal requests still use normalized residual signs and
compressed skips.

\begin{theorem}[Improved critical residual bound]
\label{thm:residualfused}
For row norms at most one, there is one uniform exact residual
sampler with the stated preprocessing budget and expected online
bit cost
\[
O\!\left(\bits^4(1+\alpha D)^{35/32}\right).
\]
Its expected number of compressed proposals and input leaves,
together with one root operation, is
\begin{equation}\label{eq:residualfused}
O((1+\alpha D)^{35/32}).
\end{equation}
The parameter $\alpha=0$ returns $x_1$ directly.
\end{theorem}
\paragraph{Proof idea.}
Expand the raw output along its last active residual update, preserving its
output gate. Above an effective-depth cutoff, the mixed potential estimate
controls each active branch; below it, the original residual sampler bounds
the remaining subtree.

\begin{proof}
Fix
\[
\eta=1/64,\quad \phi(v)=\sqrt{65/64-v^2},\quad
c_0=17/16,\quad M_0=2^{22},\quad
\kappa=3c_0=51/16,\quad J=2^{28}.
\]
For $\alpha>0$, let $K=\min\{D,\lceil J/\alpha\rceil\}$.
Use the original majority and race below or at layer $K$.
Above $K$, use the fused factory when its rational eligibility
conditions hold and the original large-bias factory otherwise.
The cutoff depends only on the weights and $\alpha$.

Consider a normalized request at layer $j+1$ above the cutoff.
Write $r=r_j$, $q_*=\tanh r$, and
$f=(1-\alpha)r+\alpha q_*$. For the actual row let
$q=\tanh\rho\le q_*$ and write
$v_0=kF(z)$ for its normalized tanh mean before radius scaling.
Put
\[
\omega=\alpha q_*/f\le\alpha,\qquad
\theta=f/r_{j+1}\le1,\qquad
w=(q/q_*)v_0,\qquad
u=(h_{j,i}-c_{j,i})/r_j.
\]
The normalized residual mean is
$v=\theta[(1-\omega)u+\omega w]+\delta_j$ with
$|\delta_j|\le\alpha\varepsilon/r_{j+1}$.
Lemma~\ref{lem:fusedlocal}, concavity, and evenness give an
expected sum of child potentials at most
\[
\theta\{(1-\omega)\phi(u)+\omega e^{\xi_j}\phi(w)\},
\qquad
\xi_j=c_0r_j^2+M_0r_j^4.
\]
Indeed, the actual-row factor $q/q_*$ can be absorbed using
$(q/q_*)\phi(v_0)\le\phi((q/q_*)v_0)$.
A zero row has no nonlinear child requests and satisfies the
same upper bound by setting $w=0$.
The ratio $\theta/\phi(\theta v)$ increases with $\theta$,
and $|\frac d{dv}\log\phi(v)|\le64$.
Thus \eqref{eq:respotentialmix} bounds the weighted offspring
factor by the exponential of
\begin{equation}\label{eq:resweightedlocal}
\alpha c_0r_j^2+
\alpha(M_0+8192)r_j^4+
64\alpha\varepsilon/r_{j+1},
\end{equation}
provided $\xi_j\le2r_j^2\le1$.

If the upper region is nonempty, then $D\ge K\ge J$.
The ideal critical radius comparison and the rounding bound give,
for $j\ge K$,
\begin{equation}\label{eq:rescutoffbounds}
r_j^2\le3/J,
\qquad
\alpha\sum_{j=K}^{D-1}r_j^4\le14/J.
\end{equation}
For the second inequality, the ideal sum is at most
$25/(4J^2)+25/(4J)$. The rounding contribution is at most
$6\varepsilon D^2<1/J$. The first follows from
$R_j^2\le(1+2\alpha j/5)^{-1}$ and
$r_j^2-R_j^2\le6D\alpha\varepsilon$.
These bounds imply $\rho\le1/32$ and
$\xi_j\le(c_0+3M_0/J)r_j^2<2r_j^2<1$.

Set $d_j=1-\tanh(r_j)/r_j$.
Since $d_j\ge r_j^2/3-r_j^4/5$,
\[
\alpha c_0r_j^2
\le\kappa\alpha d_j+(\kappa/5)\alpha r_j^4
\le\kappa\log\frac{r_j}{f_{\alpha,1}(r_j)}
       +(\kappa/5)\alpha r_j^4.
\]
Replacing $f_{\alpha,1}(r_j)$ by $r_{j+1}$ costs at most
$3\kappa\alpha\varepsilon/f_{\alpha,1}(r_j)$.
Consequently a suffix of weighted offspring factors, from level
$\ell$ to level $k\ge K$, is at most
\begin{equation}\label{eq:resweightedsuffix}
e^2(r_k/r_\ell)^\kappa.
\end{equation}
To check the constant, the fourth-order coefficient in
\eqref{eq:resweightedlocal} and the preceding display is less
than $2^{23}$. Its total is at most
$2^{23}\cdot14/J=7/16<1$.
The center and radius rounding terms total at most
\[
74\alpha D\varepsilon\sqrt{1+\alpha D}<1.
\]

Let $C_\ell(i,x)$ denote the expected number of proposals and
input leaves from one normalized request. The original sampler
below the cutoff obeys
\begin{equation}\label{eq:reslowersubtree}
C_\ell(i,x)\le7e^{21}r_\ell^{-5}\qquad(\ell\le K),
\end{equation}
by the proof of \eqref{eq:rescriticalupper} before its root gate.
The proof depends only on interval validity, so it also applies
to these midpoint centers.
Since $1/8\le\phi\le\sqrt{65}/8$, weighted charging above the
cutoff gives
\[
\frac{C_\ell(i,x)}{\phi(v_{\ell,i})}
\le e^2r_\ell^{-\kappa}
 \left(56e^{21}r_K^{\kappa-5}
       +16\alpha\sum_{k=K+1}^{\ell}r_k^\kappa\right).
\]
Here the factor $16\alpha$ charges a proposal at a reached call:
its proposal probability is at most $2\alpha$ and its potential
is at least $1/8$.
The sum obeys $\alpha\sum r_k^\kappa\le7$.
One may bound it by the third-power sum, whose ideal integral is
five, and then use the rounding estimate.
Also $r_K^{\kappa-5}\ge1$ and
\[
r_K^{\kappa-5}\le(1+\alpha K)^{29/32}
\le(J+2)^{29/32}.
\]
Together with \eqref{eq:reslowersubtree}, this proves, for every
layer including those below the cutoff,
\begin{equation}\label{eq:resuniformnormalized}
C_\ell(i,x)\le A_0\phi(v_{\ell,i})r_\ell^{-\kappa},
\qquad A_0=64e^{24}(J+2)^{29/32}.
\end{equation}

Finally condition on the raw root's last active layer $j$.
Its source radius is $r=r_{j-1}$.
Above the cutoff, Lemma~\ref{lem:fusedlocal} bounds the expected
sum of child potentials by
$r e^{\xi_{j-1}}\max\phi<4r$.
Below the cutoff the original joint bound gives at most
$r\chi_1(r)\max\phi<12r$, since
$\chi_1(r)\le1+4d_1(r)+5r^4\le10$.
Thus $16r$ is valid in both regions.
Equation \eqref{eq:resuniformnormalized} bounds the expected
descendant cost of this root by
\[
16 A_0 r^{1-\kappa}
\le16 A_0(1+\alpha D)^{35/32}.
\]
The case with no active layer costs one input probe.
Averaging over the reverse geometric choice and adding the root
operation proves \eqref{eq:residualfused}:
$J+2\le2^{29}$, $(J+2)^{29/32}<2^{27}$,
$e^{24}<2^{36}$, and $16\cdot64=2^{10}$ leave room in
the constant $2^{76}$.

The exactness induction, geometric jump implementation, and
predictable scalar-cost charging are unchanged. The fixed-precision
oracle for the fused tail has geometric index and precision tails,
so its expected scalar work per reached proposal is $O(\bits^4)$.
No extra coefficient table is stored.
\end{proof}

\begin{lemma}[A residual radius envelope above unit norm]
\label{lem:residualsuperradius}
For $s=1+\delta$ with $0\le\delta\le1$, let
$R_0=1$ and $R_{j+1}=f_{\alpha,s}(R_j)$. Then
\begin{equation}\label{eq:residualsuperradius}
R_j^2\le10\delta+\frac1{1+\alpha j/5}.
\end{equation}
\end{lemma}
\paragraph{Proof idea.}
Compare the squared radius with its supercritical equilibrium scale. Above
that scale each residual step decreases its excess by a reciprocal-square
recurrence, giving the stated effective-depth decay.

\begin{proof}
The ideal radius sequence decreases and remains in $[0,1]$.
For $0\le r\le1$, the unit Lipschitz bound and the critical
radius inequality give
\[
f_{\alpha,s}(r)\le r(1+\alpha\delta-\alpha r^2/5).
\]
If $v=r^2\ge10\delta$, the right side is at most
$r(1-\alpha v/10)$. Since
$(1-x)^2(1+2x)\le1$ for $0\le x\le1/10$,
the next squared radius is at most $v/(1+\alpha v/5)$.
Put $w_j=(R_j^2-10\delta)_+$.
If $w_j=0$, monotonicity of the radius sequence gives
$w_{j+1}=0$. Otherwise the preceding bound and direct rational
algebra give
$w_{j+1}\le w_j/(1+\alpha w_j/5)$.
Since $w_0\le1$, induction proves
$w_j\le(1+\alpha j/5)^{-1}$.
\end{proof}

\begin{proposition}[A fixed supercritical residual rate]
\label{prop:residualfusedrate}
Fix $s=1+\delta$ with $0<\delta\le2^{-16}$, let
$r_*\in(0,1)$ solve $r_*=\tanh(sr_*)$, and set
$u_*=sr_*$. There is an explicitly computable rate
$\Gamma_s=(17/16)u_*^2+O(u_*^4)$ such that the midpoint-centered
residual sampler satisfies, for every actual row norm bound $S\le s$,
\begin{equation}\label{eq:residualfusedrate}
\E J_{\mathrm f}=O_s(e^{\alpha D\Gamma_s}),\qquad
\frac{51}{16}\delta\le\Gamma_s
\le\frac{51}{16}\delta+O(\delta^2).
\end{equation}
Its expected bit cost is $O_s(\bits^4e^{\alpha D\Gamma_s})$.
The constants in the rate estimate are absolute; the prefactor is
uniform in $\alpha$. Explicit choices are given in the proof.
The sampler computes the actual reference gain and need not
receive the real comparison parameter $s$.
\end{proposition}
\paragraph{Proof idea.}
Above the cutoff, compare the weighted cost polynomial with its value at the
limiting radius. Contraction makes the excess summable, and the compressed
residual branches contribute the effective depth to the exponential rate.

\begin{proof}
Take $J=2^{28}$ and
\[
\Gamma_s=\frac{17}{16}u_*^2+2^{32}u_*^4,\qquad H_0=2^{22},
\qquad E_\delta=1+\frac{3H_0}{200}+\frac{8H_0}{\delta},
\]
\begin{equation}\label{eq:residualrateconstant}
C_\delta=2+2^{13}e^{E_\delta}
             \left(\exp(2^{26})+2+\frac{32}{51\delta}\right).
\end{equation}
Write $M_* =\exp(2^{26})$. We prove the explicit bound
$\E J_{\mathrm f}\le C_\delta e^{\alpha D\Gamma_s}$.
Let $K=\min\{D,\lceil J/\alpha\rceil\}$ for $\alpha>0$.
The case $\alpha=0$ is immediate.
By comparison with the ideal radii for gain $s$,
Lemma~\ref{lem:residualsuperradius}, and the rounding estimate,
every squared row radius above the cutoff is bounded by
\[
x_j=s^2r_j^2\le11\delta+7/J<2^{-12}.
\]
Here $s^2<11/10$, $\alpha j\ge J$, and $D\ge J$;
the radius-rounding term is less than $1/J$.
The fused factory is eligible whenever its bias is small, and
the original large-bias branch has the same weighted estimate.
For $x\le2^{-12}$, put
$\xi=(17/16)x+2^{22}x^2$.
Then $\xi\le1026x<1$, so Lemma~\ref{lem:respotentialmix}
gives a residual logarithmic offspring bound
\begin{equation}\label{eq:resratelocal}
\alpha\left(\frac{17}{16}x_j+2^{32}x_j^2\right)
+64\alpha\varepsilon/r_{j+1}.
\end{equation}
Indeed,
$2^{22}+2048\cdot1026^2<2^{32}$.
The derivative of $P(x)=(17/16)x+2^{32}x^2$ on this interval
is less than $H_0=2^{22}$.

Write $\lambda_s=s(1-r_*^2)$ and
$\lambda_\alpha=1-\alpha+\alpha\lambda_s$.
The ideal comparison radii satisfy
\[
0\le R_j-r_*\le(1-r_*)\lambda_\alpha^j.
\]
The fixed-point inequalities proved earlier give
$1-\lambda_s\ge\delta$.
Consequently the sum of the excess of \eqref{eq:resratelocal}
over $\alpha\Gamma_s$, on any suffix above the cutoff, is at
most
\[
\frac{2s^2H_0(1-r_*)}{1-\lambda_s}
+24H_0\alpha^2D^2\varepsilon
+64\alpha D\varepsilon\sqrt{1+\alpha D}
\le E_\delta.
\]
The cancellation
$\alpha/(1-\lambda_\alpha)=1/(1-\lambda_s)$ is the reason
the transient bound is uniform in $\alpha$.
Thus the weighted suffix factor from level $\ell$ to $k\ge K$
is at most $e^{E_\delta}e^{\alpha(\ell-k)\Gamma_s}$.

Below the cutoff, the original normalized residual sampler has
expected cost at most
\[
10e^{30}(J+2)^{5/2}e^{7000\delta(J+1)}\le M_*.
\]
This is the unweighted bound before its formal radius gate,
using $r_K^{-1}\le\sqrt{1+\alpha K}$.
The numerical inequality follows from $\delta\le2^{-16}$,
$J=2^{28}$, and
$7000\cdot2^{-16}(J+1)+30+\log10+(5/2)\log(J+2)<2^{26}$.

Weighted charging now gives, uniformly in the normalized output
mean at level $\ell$,
\[
C_\ell(i,x)\le
8\phi(v_{\ell,i})e^{E_\delta}e^{\alpha\ell\Gamma_s}
\left(M_*+2+\frac{2}{\Gamma_s}\right).
\]
To see that the geometric generation sum is uniform in
$\alpha$, use
\[
\frac{\alpha}{1-e^{-\alpha\Gamma_s}}
\le\alpha+\frac1{\Gamma_s}\le1+\frac1{\Gamma_s}.
\]
The proposal probability contributes its factor $2\alpha$.
For levels below $K$, the same bound follows directly from $M_*$.

At a raw root, the expected sum of child potentials is at most
1024. Below the cutoff this follows from
$\chi_g(r)\le10e^4<600$ and $\max\phi<9/8$;
above it the weighted bound and $\xi<1$ give a smaller constant.
Conditioning on the reverse geometric choice of the last active
layer, using $\Gamma_s\ge51\delta/16$, and adding the root
operation proves \eqref{eq:residualrateconstant} and
the work bound.

Finally, $3\delta\le u_*^2\le3\delta+(3/2)\delta^2$ gives
\[
\Gamma_s\le\frac{51}{16}\delta+
 \left(\frac{51}{32}+2^{32}(25/8)^2\right)\delta^2
\le\frac{51}{16}\delta+2^{36}\delta^2.
\]
Its lower bound follows from $u_*^2\ge3\delta$.
The finite-bit implementation and exactness induction are those
of Theorem~\ref{thm:residualfused}.
\end{proof}
